\scanpage{26}
\textbf{C. R. Acad. Sc. Paris, t. 279 (22 juillet 1974)}\hfill\textbf{Série A ---}\mspage{129}\par\noindent\rule{\linewidth}{.35pt}\par

TOPOLOGIE. --- \textit{Classification topologique des plongements des 1-complexes compacts dans les surfaces.} Note (*) de M. \textbf{Yves Ladegaillerie}, présentée par M. Henri Cartan.

{\small On établit l'équivalence de la catégorie isotopique des voisinages tubulaires orientés des 1-complexes avec la catégorie des 1-complexes circularisés. On établit ensuite l'équivalence de la catégorie isotopique des voisinages tubulaires quelconques des 1-complexes avec la catégorie des 1-complexes munis d'un automorphisme involutif sans point fixe et d'une circularisation antisymétrique. Enfin on caractérise les plongements dans des surfaces orientées de genre minimal.}

7. \textsc{Voisinages tubulaires des 1-complexes.} --- La terminologie et les notations sont celles de (\textsuperscript{1}). Dans la suite \(\mathrm{K}\) désigne un 1-complexe topologique compact.

7.1. \textsc{Théorème.} --- \textit{Soit \(i\) un plongement de \(\mathrm{K}\) connexe dans une surface à bord compacte orientée \(\mathrm{X}\), soit \(\mathrm{A}\) un anneau non réduit à \(0\). Il y a équivalence entre :}

\textit{(a) \(i\) est une équivalence d'homotopie ;}

\textit{(b) \(i\) est équivalent à un plongement standard \(i_\omega\) [cf.} (\textsuperscript{1}), 3.3] ;

\textit{(c) \(x\mathrel{\ooalign{\hfil$\simeq$\hfil\cr\hfil$\vert$\hfil\cr}}\mathrm{S}^2\), \(\pi_0(i)\) \textup{et} \(\pi_1(i)\) sont des isomorphismes ;}

\textit{(d) \(\mathrm{X}\mathrel{\ooalign{\hfil$\simeq$\hfil\cr\hfil$\vert$\hfil\cr}}\mathrm{S}^2\), \(\mathrm{X}\) connexe et \(\mathrm{H}_1(i)\) est un isomorphisme de \(\mathrm{H}_1(\mathrm{K}\,;\mathrm{A})\) sur \(\mathrm{H}_1(\mathrm{X}\,;\mathrm{A})\) ;}

\textit{(e) \(\mathrm{X}\mathrel{\ooalign{\hfil$\simeq$\hfil\cr\hfil$\vert$\hfil\cr}}\mathrm{S}^2\), \(\mathrm{X}\) connexe et \(\mathrm{H}^1(i)\) est un isomorphisme de \(\mathrm{H}^1(\mathrm{K}\,;\mathrm{A})\) sur \(\mathrm{H}^1(\mathrm{X}\,;\mathrm{A})\).}

On dira alors que \((\mathrm{X},i)\) est un \textit{voisinage tubulaire orienté de} \(\mathrm{K}\).

7.2. Plus généralement on appelle \textit{voisinage tubulaire} d'un 1-complexe \(\mathrm{K}\) tout couple \((\mathrm{X},i)\) où \(i\) est un plongement de \(\mathrm{K}\) dans la surface compacte à bord \(\mathrm{X}\) et \(i\) est une équivalence d'homotopie. Ce sont les objets de la \textit{catégorie isotopique des voisinages tubulaires} notée TUB, un \textit{morphisme entre les voisinages tubulaires} \((\mathrm{X},i)\) et \((\mathrm{X}',i')\) de \(\mathrm{K}\) et \(\mathrm{K}'\) étant une classe d'isotopie d'homéomorphismes \(h\) de \(\mathrm{X}\) sur \(\mathrm{X}'\) tel\unsure{\textperiodcentered|.} que \(h(i(\mathrm{K}))=i'(\mathrm{K}')\).

7.3. La \textit{catégorie isotopique des voisinages tubulaires orientés} est la catégorie TUBO dont les objets sont les voisinages tubulaires orientés de 1-complexes et dont les morphismes sont les morphismes de voisinages tubulaires qui respectent l'orientation.

8. \textsc{La catégorie isotopique des 1-complexes circularisés.}

8.1. C'est la catégorie \(\mathrm{K}\,\omega_{\mathrm{I}}\) dont les objets sont les couples \((\mathrm{K},\omega)\) où \(\mathrm{K}\) est un 1-complexe quelconque et \(\omega\) une circulation de \(\mathrm{K}\). Un morphisme de \((\mathrm{K},\omega)\) dans \((\mathrm{K}',\omega')\) est une classe d'isotopie d'homéomorphismes de \(\mathrm{K}\) dans \(\mathrm{K}'\) qui respecte la circulation.

8.2. On a un foncteur \(\mathrm{T}\) de \(\mathrm{K}\,\omega_{\mathrm{I}}\) dans TUBO qui à \((\mathrm{K},\omega)\) associe le voisinage tubulaire standard \(\mathrm{T}(\mathrm{K},\omega)\) (\textsuperscript{1}) et qui à un morphisme \(\bar h:(\mathrm{K},\omega)\to(\mathrm{K}',\omega')\) associe le morphisme \(\mathrm{T}(\bar h)\) de voisinages tubulaires classe de l'homéomorphisme \(\mathrm{H}:\mathrm{T}(\mathrm{K},\omega)\to\mathrm{T}(\mathrm{K}',\omega')\) défini à partir des homéomorphismes \(\mathrm{H}_a:\bar a\times\mathrm{I}\to\unsure{\overline{h(a})|\overline{h(a)}}\times\mathrm{I}\) (\(a\) étant un arc de \(\mathrm{K}\)) qui à \((x,t)\) associe \((h(x),t)\).

% right margin of 8.2, a shallow curved pen stroke and two dots
\figure{\begin{tikzpicture}[x=1cm,y=1cm]\draw[thin] (0,0) .. controls (.44,-.12) and (.93,-.18) .. (1.05,.04);\fill (.32,.37) circle (.012);\fill (.39,.44) circle (.008);\end{tikzpicture}}

8.3. \textsc{Théorème.} --- \textit{\(\mathrm{T}\) est une équivalence entre les catégories \(\mathrm{K}\,\omega_{\mathrm{I}}\) et TUBO.}

9. \textsc{La catégorie isotopique des \(\mathbf{Z}_2\)-voisinages tubulaires orientés.}

9.1. C'est la catégorie \(\widetilde{\mathrm{TUBO}}\) dont les objets sont les couples \((\widetilde{\mathrm{T}},\theta)\) où \(\widetilde{\mathrm{T}}\) est un objet de TUBO et \(\theta\) un homéomorphisme de voisinage tubulaire de \(\widetilde{\mathrm{T}}\) sur \(\widetilde{\mathrm{T}}\), involutif, sans
% archive sheet number: 25

